\section{Пять осей персонажа: аффект, социальная роль, дистанция/близость, способ заботы, способ конфликта}

Если типаж — это набор правил, а характер — правила плюс причина, то нужен язык, чтобы описывать сами правила системно, а не интуитивно ("она как бы холодная, но по-своему тёплая"). Удобно представить персонажа не как список прилагательных, а как точку в пространстве нескольких осей — и здесь мы возьмём пять базовых, оставив полную девятиосевую версию для главы 5.

Аффект — эмоциональный регистр по умолчанию: тёплый, холодный, взрывной, ровный. Социальная роль — позиция относительно других: лидер, помощница, соперница, наблюдательница. Дистанция/близость — на каком расстоянии персонаж комфортно держит других и как быстро подпускает ближе. Способ заботы — что именно персонаж делает, когда ему не всё равно: кормит, чинит, ругает, защищает, шутит, дразнит. Способ конфликта — что персонаж делает, когда сталкивается с сопротивлением: спорит, обижается, исчезает, атакует, рационализирует.

Пять осей — не анкета для заполнения, а система координат, которая делает типажи сравнимыми друг с другом не по названию, а по факту. Два персонажа с формально разными ярлыками ("kuudere" и "genki girl") могут оказаться рядом по трём осям из пяти и различаться только по одной — и это будет куда точнее, чем спорить, какой ярлык вернее. И наоборот: два персонажа с одинаковым ярлыком tsundere могут разойтись по осям дистанции и способа конфликта настолько, что окажутся дальше друг от друга, чем от совершенно другого жанрового типа.

Тест на архетип-как-инструмент из 0.5 теперь можно сформулировать точнее: собрать Рин "из тропов" значит выбирать значения по каждой оси по одному, без связи между ними. Вырастить характер значит убедиться, что значение на одной оси причинно объясняет значения на других, — что её способ заботы (резкая, техническая) логически следует из её дистанции (близкая, но не сливающаяся) и её способа конфликта (спорит, но не бросает), а не подобраны отдельно, потому что "так эффектнее".
